\documentclass{article}
\usepackage{hyperref}
\usepackage{Style}

\nocite{*} % Comentar si quiero citar
%\addbibresource{bibliografia.bib} % Quitar el comentado si quiero usar bibliografia

\begin{document}

\begin{minipage}{2.5cm}
    \includegraphics[width=2cm]{imagen_puc.jpg}
\end{minipage}
\begin{minipage}{12cm}
    {\sc Pontificia Universidad Católica de Chile\\
    Facultad de Matemáticas\\
    Departamento de Matemática\\
    Profesor: Manuel Cabezas -- Ayudante: Benjamín Mateluna}
\end{minipage}
\vspace{2ex}

{\centerline{\bf Introducción a la Geometría - MAT1304}
\centerline{\bf Ayudantía 4 - Repaso I1}}
\centerline{\bf 19 de marzo de 2026}

\vspace{1cm}
\noindent\textbf{Problema 1.} Sea $\triangle ABC$ isósceles de base $\overline{BC}$, sean $E$ y 
$D$ puntos en $\overline{AB}$ y $\overline{AC}$ respectivamente, de forma que 
$\overline{AE}=\overline{AD}$. Demuestre lo siguiente:
\begin{enumerate}
    \item Demuestre que $\overline{ED}\parallel\overline{BC}$.
    \item Sea $P\in\overline{EC}\cap\overline{DB}$, demuestre que $\overline{AP}$ es bisectriz de
    $\overline{BAC}$.
\end{enumerate}

\vspace{5mm}
\noindent\textbf{Problema 2.} En $\triangle ABC$, un punto $D$ esta sobre $\overline{AC}$ tal que
$\overline{AB}=\overline{AD}$. Asumir que $\angle ABC-\angle ACB=30^{\circ}$. Encontrar 
$\angle CBD$.

\vspace{5mm}
\noindent\textbf{Problema 3.} Sea $\triangle ABC$ tal que $\angle ABC=\angle BCA=40^{\circ}$ y sea
$D$ el punto en el lado $\overline{AC}$ de modo que $\angle ABD=\angle DBC=20^{\circ}$. Demuestre 
que $\overline{AD}+\overline{BD}=\overline{BC}$.

\vspace{5mm}
\noindent\textbf{Problema 4.} Se construyen exteriormente a un $\triangle ABC$ tres triángulos 
equiláteros $\triangle ABD$, $\triangle BCF$ y $\triangle CAE$. Muestre que 
$\overline{AF}=\overline{BE}=\overline{CD}$.

\vspace{1cm}
\noindent\textbf{\underline{Ejercicios Propuestos:}}

\vspace{5mm}
\noindent\textbf{Extra 1.} Sean $\triangle ABE$, $\triangle BED$ y $\triangle BDC$ los 
correspondientes triángulos incluyendo su interior, tales que:
\begin{itemize}
\item $\triangle ABE\cap\triangle BED=\overline{BE}$, $\triangle ABE\cap\triangle BDC=B$.
\item $\triangle BED\cap\triangle BDC=\overline{BD}$.
\item $\overline{AE}=\overline{BE}$ y $\overline{BD}=\overline{CD}$.
\end{itemize}
Dibuje la situación y demuestre las siguientes afirmaciones.
\begin{enumerate}
    \item Si $A,B,C$ son colineales, es decir, están contenidos en una misma recta, entonces las
    rectas $\overleftrightarrow{AE}$ y $\overleftrightarrow{CD}$ son secantes.
    \item Las rectas $\overleftrightarrow{AE}$ y $\overleftrightarrow{CD}$ son paralelas si y solo
    si
    \begin{equation*}
        2\angle ABE+2\angle DBC+\angle EBD=360^{\circ}
    \end{equation*}
\end{enumerate}

\vspace{5mm}
\noindent\textbf{Extra 2.} Cuatro cuadrados son construidos exteriormente a los lados de un 
paralelogramo. Mostrar que los centros de los cuadrados forman un cuadrado.

%\printbibliography % Quitar el comentado si quiero usar bibliografia

\end{document}
